\section{Conclusion}

We introduce paired human and mouse benchmarks for tissue-specific MHC-I presentation preference. Standard evaluations establish task learnability, while matched peptide-disjoint OOF reveals a consistent unseen-peptide generalization gap. On identical strict folds, human auxiliary supervision retains a reproducible benefit, but the TissuePMHC human model is comparable to the strongest auxiliary dual-branch control; the mouse Factorized MMoE is likewise comparable to shared heads and a BLOSUM62 random forest, with much of its final performance attributable to seed averaging. Capacity-matched MHC-only controls and frozen MHCflurry/NetMHCpan controls show that general peptide--MHC signal explains part, but not all, of the tissue-conditioned performance. The main contribution is therefore a reproducible framework for separating task learnability, unseen-peptide robustness, architecture effects, and general peptide--MHC signal without claiming universal architecture, unseen-task, external-cohort, or clinical generalization.
